% !Mode:: "TeX:UTF-8"

\chapter{相关研究综述}[Related Work Review]

在第 1 章明确了本文关注的人形机器人运动技能学习问题之后，本章进一步从相关研究层面展开梳理。需要回顾人形机器人敏捷控制、全身遥操作、大模型驱动具身智能以及人体动作模仿等研究进展，以明确本文工作的技术脉络和切入点，并为后续方法章节提供清晰的研究背景支撑。

\section{人形机器人敏捷控制与具身智能研究进展}[Agile humanoid control and embodied intelligence]

人形机器人全身运动控制的核心目标，是在满足动力学约束和接触稳定性的前提下，使机器人完成自然、协调且具有任务适应性的运动。早期研究主要集中在基础移动控制任务，例如平地行走、速度跟踪和简单避障。在这类任务中，下肢通常被视为实现移动的主要机构，控制策略更关注质心稳定和足端接触，而对上肢、躯干和全身姿态表达能力的利用相对有限。

近年来，随着高性能人形机器人硬件和并行物理仿真平台的发展，研究者开始关注更复杂的全身运动技能。基于强化学习的方法已经在双足行走、快速移动、跳跃和复杂地形运动中取得进展\cite{xie2020learning,li2023robust,radosavovic2024real,li2024reinforcement,radosavovic2402humanoid,zhuang2024humanoid,gu2024advancing,long2024learning}。这类方法通常在仿真环境中进行大量试错，通过奖励函数引导策略形成稳定运动模式。与传统控制方法相比，学习型控制能够在复杂接触和高维状态空间下自动搜索控制规律，减少对人工设计步态机和接触序列的依赖。

除基础移动能力外，整身控制和表达性动作逐渐成为人形机器人研究的重要方向。相关工作尝试让机器人完成舞蹈、上肢跟随、全身远程操作、移动操作和高动态运动等任务\cite{cheng2024expressive,he2024learning,he2024omnih2o,zhang2024wococo,fu2024humanplus,ji2024exbody2}。这类任务要求策略同时协调下肢支撑、躯干姿态和上肢运动，不仅要保持动态平衡，还要尽可能复现参考动作的节奏、幅度和姿态风格。因此，人形机器人控制问题正在从“稳定行走”扩展到“自然、敏捷、可表达的全身运动”。其中，全身遥操作任务通常需要同时处理人体测量信息、动作重定向、机器人局部反馈以及环境感知反馈，其典型闭环架构如图~\ref{fig:agile_humanoid_research} 所示。

与普通移动控制相比，敏捷全身运动具有更高难度。首先，机器人自由度高，状态空间和动作空间维度大，简单的人工控制规则难以覆盖全部运动模式。其次，足端接触状态频繁变化，腾空、落地和支撑切换会引入强非线性冲击。再次，许多动作具有明显的时序结构，某一时刻的控制误差会累积到后续相位，导致整段动作失败。因此，本文将动作表示、参考轨迹构建和强化学习闭环控制作为一个整体问题进行研究。

\begin{figure}[htbp]
  \centering
  \vspace{2mm}
  \includegraphics[width=0.98\textwidth]{遥操作.png}
  \vspace{2mm}
  \caption{人形机器人全身遥操作任务的典型系统架构示意图}
  \label{fig:agile_humanoid_research}
\end{figure}

随着具身智能研究的持续升温，大模型正逐渐成为连接机器人感知、认知、规划与控制的重要技术底座。相关综述指出，大语言模型、视觉 Transformer、视觉语言模型、视觉生成模型以及具身多模态大模型，已经开始在环境理解、任务推理、语言交互、动作生成和行为控制等环节发挥作用\cite{zeng2025embodied}。对于人形机器人而言，这类模型的意义在于：一方面，它们能够利用互联网规模数据获得较强的语义先验和跨模态关联能力；另一方面，它们可以把开放环境中的图像、文本、语音和视频信息统一到同一决策框架中，从而提升机器人在非结构化场景下的任务适应性。

\begin{figure}[htbp]
  \centering
  \vspace{2mm}
  \includegraphics[width=0.98\textwidth]{大模型具身智能.png}
  \vspace{2mm}
  \caption{一种基于分布式多模态大模型的人形机器人作业技术路线图}
  \vspace{4mm}
  \label{fig:distributed_multimodal_llm_humanoid}
\end{figure}

现有大模型驱动的人形机器人研究大体可归纳为三类技术路线。第一类是分布式模块化路线，即分别使用不同模型处理感知、规划、决策和控制等子任务，再通过接口或中间表示进行级联。这一路线的优点是结构清晰、可解释性较强，也便于将现有视觉模型和语言模型快速接入机器人系统；其局限在于模块间误差容易累积，推理延迟和接口对齐问题会影响整机闭环性能。第二类是端到端一体化路线，即直接从视觉和语言输入映射到动作输出，尝试利用视觉--语言--动作模型统一高层语义与低层控制。该路线具备较强的任务泛化潜力，但通常依赖大规模机器人交互数据和高质量标注，且在高频、强接触的人形机器人全身运动控制上仍面临训练成本高和稳定性不足的问题。第三类是云边端协同路线，即将高算力语义推理、长期知识存储和低时延执行控制分散到云端、边缘端和本体端协同完成，以平衡算力、隐私和实时性需求\cite{zeng2025embodied}。

其中，分布式多模态大模型路线的一个典型思路，是围绕传感器、控制器和执行器构建分层的感知、认知、决策与控制模型，并通过迭代训练方式实现任务语义理解到执行控制的逐级映射。其代表性技术框架如图~\ref{fig:distributed_multimodal_llm_humanoid} 所示。

从本文研究主题来看，大模型研究进展为人形机器人动作学习提供了重要启发，但并未直接解决敏捷全身运动中的核心困难。大模型更擅长回答“做什么”和“为什么这样做”，而本文关注的动作跟踪与技能学习更强调“如何在动力学约束下稳定做到”。特别是在跳跃、快速转向、舞蹈等高动态动作中，机器人必须同时满足关节限位、接触约束、姿态协调和时序连续性要求，因此仍需要依赖高质量的人体动作先验、合理的动作重定向方式以及物理仿真中的强化学习闭环控制。也就是说，大模型可以为人形机器人提供更强的高层语义条件和任务上下文，而稳定的底层运动技能学习框架则是这些高层能力落地的基础。

\section{基于强化学习的人体动作模仿与数据驱动范式}[Reinforcement Learning based motion imitation and human data driven paradigm]

在物理角色动画领域，基于强化学习的动作模仿已经被证明能够生成自然且具有动态一致性的运动。DeepMimic 通过设计参考姿态、速度、末端位置和质心等跟踪奖励，使物理角色能够学习翻滚、踢腿和后空翻等复杂技能\cite{deepmimic2018}。随后，对抗运动先验（Adversarial Motion Priors, AMP）、对抗技能嵌入（Adversarial Skill Embeddings, ASE）、MaskedMimic 和 Perpetual Humanoid Control 等方法进一步研究了风格化运动先验、可复用技能表示、部分观测条件下的动作补全和长时序角色控制\cite{peng2021amp,peng2022ase,tessler2024maskedmimic,luo2023perpetual}。这些研究说明，参考动作不仅可以作为监督目标，还可以作为强化学习中的运动先验，约束策略在物理环境中探索合理动作。

将上述思想应用到机器人控制时，需要进一步考虑机器人结构和执行器约束。机器人不是理想化物理角色，其关节自由度、关节限位、执行器带宽和足端接触模型均会限制动作可行性。已有研究开始将人体参考动作引入人形机器人控制，例如利用人体动作进行全身远程操作、整身模仿或带有人体参考的运动控制\cite{he2024learning,zhang2024whole,fu2024humanplus}。这类方法的共同特点是先构建适合机器人执行的参考轨迹，再通过强化学习或策略优化获得闭环控制器。

本文采用的技术路线与上述研究具有一致思想：首先从人体运动数据中提取全身运动先验，再将其转化为机器人可执行的参考轨迹，最后在物理仿真环境中训练策略进行跟踪。与仅学习行走速度指令的任务相比，动作模仿任务需要同时关注根节点运动、关节姿态、刚体位置、速度和动作相位，因此对观测设计、奖励函数和初始化策略提出了更高要求。

为了让机器人获得复杂且自然的运动技能，本文采用基于真实人体运动数据的模仿学习思路。人体运动捕捉数据通常包含丰富的动作类别和连贯的全身姿态变化，能够为机器人提供高质量的运动先验。典型流程包括人体运动数据解析、动作重定向、参考轨迹生成和基于强化学习的动作跟踪\cite{he2024learning,tessler2024maskedmimic}。

在数据层面，AMASS 等大规模人体运动捕捉数据集将不同来源的人体动作统一表示为 SMPL 参数化身体模型\cite{amass2019}。该表示能够提供人体根位姿、关节旋转和身体形状参数，便于进行跨数据集处理。但由于人体与机器人在骨骼拓扑、关节轴定义和连杆尺寸上存在差异，需要通过动作重定向将人体关键点运动映射到机器人关节空间。

动作重定向是连接人体数据和机器人控制的关键环节。常见做法不是逐关节复制人体角度，而是建立人体关键点与机器人连杆之间的空间对应关系，通过优化使机器人关键点尽可能贴合人体关键点轨迹。这样可以在一定程度上绕开人体骨骼和机器人关节定义不一致的问题。与此同时，还需要考虑机器人关节上下限、根节点高度、足端接触和轨迹平滑等约束，否则生成的参考动作虽然在几何上接近人体姿态，却可能在物理仿真中难以稳定跟踪\cite{he2024learning,zhang2024whole}。

在训练层面，基于相位的强化学习动作跟踪是一种常见范式。策略网络不仅接收机器人自身的关节状态、速度和姿态信息，还接收动作相位和参考状态误差，使策略能够感知当前应当执行的动作片段。为提升训练稳定性，通常还需要引入非对称 Actor-Critic 架构、参考状态初始化、课程学习和动作偏离终止机制\cite{deepmimic2018}。此外，生成对抗模仿学习等方法也为从专家数据中学习策略提供了重要思路\cite{gail2016}。本文重点采用显式轨迹跟踪奖励，而不是对抗判别器奖励，原因在于显式误差项更便于结合刚体位置、关节角和速度等可解释指标进行调试。

\section{仿真训练中的模型误差与鲁棒性方法}[Model mismatch and robustness]

人形机器人强化学习通常依赖物理仿真器完成大规模采样。MuJoCo、Isaac Gym 和 Genesis 等平台能够提供刚体动力学、接触求解和并行训练能力，为复杂机器人控制提供了重要工具\cite{todorov2012mujoco,isaacgym2021,Genesis}。但仿真模型仍会受到连杆惯量、摩擦参数、接触模型和执行器响应等误差影响。即使本文主要关注仿真环境中的动作学习，这些模型误差也会影响策略训练稳定性和跨仿真器复现能力。

现有研究中常见的鲁棒性方法包括系统辨识、域随机化和残差学习。系统辨识通过实验数据估计质量、阻尼、摩擦和电机响应等参数，以获得更准确的仿真模型\cite{hwangbo2019learning,yu2019sim,du2021auto}；域随机化则在训练时随机扰动物理参数，使策略适应一定范围内的模型变化\cite{tobin2017domain,tan2018sim,peng2018sim}；残差学习通过额外模型补偿基础控制器或动力学模型的不足\cite{silver2018residual,johannink2019residual,alakuijala2021residual}。对于本文的研究范围而言，这些方法主要作为提升策略稳定性和泛化性的参考，而不是本文工作的核心内容。本文更关注如何先在仿真中形成稳定、可复现的动作跟踪流程。

\section{本章小结}[Summary]

本章分析了人形机器人敏捷运动控制、基于强化学习的动作模仿、人体运动数据重定向以及仿真训练鲁棒性等相关研究。后续章节将围绕专家数据预处理、强化学习动作跟踪方法以及实验验证三个部分展开。
